\section{La constante holográfica de acoplamiento aritmético-geométrico}
\label{lib01:acoplamiento}

El registro dodecafásico enlaza dos realizaciones del mismo estado: la
lectura posicional que interviene en la clausura aritmética y la lectura
que conserva incidencia, orientación y dirección. El desarrollo siguiente
determina las operaciones que permite realizar el registro ya producido,
la información que retienen sus memorias y los errores bajo los cuales
pueden recuperarse sus relaciones.
Su función de referencia admite una formulación precisa: la órbita de fase
del registro es una base completa de su portador, y la memoria isométrica
conserva exactamente el Gram de esa base.

El punto de partida sigue siendo la composición
\[
 \APP\longrightarrow\mathrm{TRIT}\longrightarrow\TPK
 \longrightarrow\text{estado enriquecido}
 \longrightarrow\text{estructura discreta conjunta del continuo}.
\]
APP aporta las hojas y sus residuos y cocientes; TRIT conserva régimen y
orientación; TPK selecciona, transporta y actualiza conservando acarreo,
ruta, frontera, supervivencia y memoria. Las cinco construcciones del
continuo pertenecen a esa producción conjunta, desarrollada en el cuerpo
precedente. Aquí se trabaja después de ella, sobre su publicación entera
dodecafásica y sus lectores tipados. Ni el portador lineal posterior ni el
conjunto de candidatos de un decodificador reemplazan el dominio de
historias generadas.

\subsection{Objeto, carta y producción del registro}
\label{lib01:registro}

La definición~\ref{lib01:definicion} fija el objeto generado y sus lectores. Se desarrolla ahora su codificación reversible, su capacidad de identificación operatoria y la estabilidad de sus realizaciones incidenciales.

La denominación expresa una función constitutiva del objeto, no una
identificación automática con un coeficiente de interacción física. La
canonicidad se refiere al perfil, origen y orientación declarados. Una
transformación de carta y una actualización que incorpora nuevos eventos
no son la misma operación. Tampoco se afirma que el mismo vector numérico
deba ser un invariante necesario de todo sistema.

El calendario de ciento ocho posiciones consta de doce ventanas nonádicas.
La extensión
\[
 0\longrightarrow C_{12}\longrightarrow C_{108}
 \longrightarrow C_9\longrightarrow0
\]
conserva el acarreo de fase; su no escisión y la distinción entre retorno de
fase y retorno del estado están demostradas en
\cref{prop:k-extension}. El registro de una historia retiene, en cada
ventana, la subruta, su orientación, el acarreo y el registro de incidencias.
El orden causal de la extracción es
\begin{equation}
 x_{\rm term}\longmapsto\mathcal R_{12}(x_{\rm term})
 \longmapsto(b^{90},b^{120},Q_{\rm TPK})
 \longmapsto U_{\rm sgn}\longmapsto K.
 \label{lib01:cadena-registro}
\end{equation}
Las coordenadas regionales y el registro firmado son dos lecturas del
mismo terminal; no se postula que la matriz ternaria visible determine por
sí sola la coordenada firmada.

Para fijar materialmente la última composición, los canales producidos son
\begin{align*}
 b^{90}&=(-495,-116,-684,108,601,-90,
                -173,-88,313,560,-397,461),\\
 b^{120}&=(-425,-281,-481,671,-255,30,
                 475,-543,680,251,6,-128),\qquad Q_{\rm TPK}=6263.
\end{align*}
Sea \(B_r=\sum_{j=0}^3b^{120}_{r+3j}\) y
\(d_{r,j}=b^{90}_{r+3j}\), para \(r=1,2,3\). El lector armónico da
\[
 s_1=\frac{Q_{\rm TPK}+2B_1+B_2}{3},\qquad
 s_2=s_1-B_1,\qquad s_3=s_2-B_2,
\]
\[
 U^{(r)}=(s_r,d_{r,1}-d_{r,3},d_{r,0}+d_{r,2},d_{r,0}-d_{r,2}).
\]
En esta salida compatible, \((B_1,B_2,B_3)=(972,-1073,101)\),
\((s_1,s_2,s_3)=(2378,1406,2479)\), y
\begin{align*}
 U^{(1)}&=(2378,-452,-668,-322),\\
 U^{(2)}&=(1406,998,-204,-28),\\
 U^{(3)}&=(2479,-551,-371,-997).
\end{align*}
La división por tres pertenece a las tres órbitas del lector y la división
por cuatro que sigue a la inversa de Hadamard; no son coeficientes ajustados
a una salida física. Sobre las órbitas
\((1,4,7,10)\), \((2,5,8,11)\), \((3,6,9,12)\), se aplica
\[
 H_4=\begin{pmatrix}
 1&1&1&1\\1&1&-1&-1\\1&-1&1&-1\\1&-1&-1&1
 \end{pmatrix},\qquad H_4^2=4I.
\]
Los productos \(H_4U^{(r)}\) son, respectivamente,
\[
 (936,2916,2484,3176),\quad(2172,2636,232,584),
 \quad(560,3296,3656,2404).
\]
Su divisibilidad por cuatro verifica la pertenencia a la subred de
Hadamard. Tras dividir y restituir el orden natural se obtiene
\begin{equation}
 K=(234,543,140,729,659,824,621,58,914,794,146,601),
 \qquad\mathbf1^*K=6263.
 \label{lib01:K}
\end{equation}
La inversión y sus congruencias, incluido el índice \(4096\) de la imagen
firmada, están probadas en \cref{lem:k-hadamard,prop:k-sello}. La
reconstrucción es única porque \(H_4^{-1}=H_4/4\), no porque se haya
buscado un registro próximo a un valor conocido de alfa.

Para evitar una colisión de símbolos, escribiremos
\(\Delta_j=I-S^j\) para los extractores enteros antecedentes, con
\((Sx)_i=x_{i+1}\) e índices módulo doce. Entonces
\[
 \Delta_3K=b^{90},\qquad \Delta_4K=b^{120},\qquad Q(K)=6263.
\]
El extractor \((\Delta_3,\Delta_4,Q)\) es inyectivo, con factores
invariantes \(1,\ldots,1,12\), según \cref{prop:k-transversal}. Sus
condiciones de imagen integral y su ley de acumulación permanecen las del
registro antecedente. Los complementos normalizados de memoria que se
denotarán por \(D_j\) son otros operadores, derivados de esos
desplazamientos y con un codominio métrico declarado.

\subsection{Codificación periódica, clausura e inversión exacta}
\label{lib01:kappa}

Fijemos \(B=1000\), doce posiciones y
\[
 N(k)=\sum_{j=1}^{12}k_jB^{12-j},\qquad d_B=B^{12}-1,
 \qquad \kappa(k)=\frac{N(k)}{d_B}.
\]
En la carta canónica,
\begin{equation}
 \begin{aligned}
 \kappa_{\rm ag}
 &=\frac{234543140729659824621058914794146601}{10^{36}-1},\\
 \alpha_A&=A_{reg}-\kappa_{\rm ag},\qquad
 A_{reg}=\pi_{\rm HMT}+e_{\rm HMT}-\varphi_{\rm HMT}-4.
 \end{aligned}
 \label{lib01:clausura}
\end{equation}
Las tres coordenadas regionales se publican antes de actuar en este lector;
su origen común con el cierre dodecafásico no obliga a confundir el orden de
evaluación de las coordenadas correlacionadas. La igualdad es un lector
posterior del estado generado. Despejarla una vez conocidas sus salidas es
una operación inversa legítima, distinta de utilizar un valor objetivo para
producir el registro.

\begin{proposition}[Reversibilidad y radio de la codificación]
\label{lib01:inversion-kappa}
En el dominio \(\{0,\ldots,999\}^{12}\), la fracción \(\kappa(k)\),
con base y longitud declaradas, determina exactamente \(k\). Lo mismo
ocurre con el par exacto \((A_{reg},\alpha_A)\) cuando satisface
\eqref{lib01:clausura}. Si
\[
 |\widetilde\kappa-\kappa(k)|<\frac1{2d_B},
\]
el redondeo de \(d_B\widetilde\kappa\) al entero más próximo recupera
\(N(k)\) sin error. Si se aproximan por separado los dos términos de la
clausura, basta que la suma de sus errores absolutos sea menor que
\(1/(2d_B)\).
\end{proposition}
\begin{proof}
La multiplicación por \(d_B\) recupera el entero \(N(k)\).
Doce divisiones euclídeas por \(B\), conservando los restos y el número
de posiciones, recuperan sus bloques, incluidos los iniciales nulos. La
hipótesis de error sitúa el entero verdadero a distancia estrictamente
menor que media unidad. Para la clausura se aplica la desigualdad
triangular al error de la diferencia. Dos enteros consecutivos muestran
que el radio no puede ampliarse uniformemente en todo el dominio del
código: su punto medio es ambiguo.
\end{proof}

Si se recibe una fracción reducida \(p/q\), se calcula
\(N=d_Bp/q\) y se comprueban integralidad y \(0\leq N\leq d_B\).
El denominador reducido no sustituye los metadatos de base, longitud y
carta. El código contiene \(10^{36}\) palabras y necesita
\(\lceil\log_2(10^{36})\rceil=120\) bits en una representación binaria
de longitud fija. Doce bloques de diez bits ocupan también 120 bits: no se
ha demostrado una compresión adicional, autenticación ni integridad
criptográfica. Una fracción exacta o el entero \(N\) evita perder la
precisión necesaria para esta inversión.

La lectura periódica no se confunde con el truncamiento \(N/B^{12}\).
Si \(\kappa_j=\kappa(S^jK)\), una división posicional da
\begin{equation}
 \kappa_{j+1}=1000\kappa_j-K_{j+1},\qquad
 K_{j+1}=1000\kappa_j-\kappa_{j+1},\qquad
 \sum_{j=0}^{11}\kappa_j=\frac{6263}{999}.
 \label{lib01:orbita-kappa}
\end{equation}
Así, la órbita escalar conserva todas las diferencias de fase y permite
leer \(B_K=\{j+1:1000\kappa_j-\kappa_{j+1}\geq729\}\).
La reversibilidad de estas publicaciones no implica que un escalar
reconstruya toda la historia enriquecida que produjo el registro.

\subsection{La órbita completa como referencia de su portador}
\label{lib01:orbita}

En \(V=\mathbb R^{12}\), con su producto euclídeo canónico, consideremos
\[
 O_K=(K,SK,\ldots,S^{11}K):\mathbb R^{12}\longrightarrow V.
\]
El desplazamiento \(S\) actúa sobre posiciones del registro. No se
identifica por la sola periodicidad con la evolución completa TPK o con
una simetría excepcional posterior.

\begin{theorem}[Referencia cíclica completa y Gram exacto]
\label{lib01:marco-ciclico}
La matriz \(O_K\) es invertible. Su Gram tiene los valores propios
\[
\begin{array}{c|c}
 \text{valor}&\text{multiplicidad}\\\hline
 39225169=6263^2&1\\
 1248025-345420\sqrt3&2\\
 589609&2\\
 1407529&2\\
 1053157&2\\
 1248025+345420\sqrt3&2\\
 697225&1
\end{array}
\]
y satisface
\begin{equation}
 589609\|c\|^2\leq\|O_Kc\|^2\leq39225169\|c\|^2,
 \quad
 \|O_K^{-1}\|=\frac1{\sqrt{589609}},\quad
 \operatorname{cond}(O_K)=\frac{6263}{\sqrt{589609}}.
 \label{lib01:cotas-orbita}
\end{equation}
La órbita del registro centrado tiene rango once.
\end{theorem}
\begin{proof}
Las entradas del Gram son las autocorrelaciones
\(\langle K,S^jK\rangle\). En orden \(j=0,\ldots,11\) valen
\begin{align*}
 (&4251257,2999323,3163385,3287920,3216553,3344743,\\
  &2950064,3344743,3216553,3287920,3163385,2999323).
\end{align*}
Como \(S\) es ortogonal, el Gram es circulante. Sustituir los doce
caracteres \(\exp(2\pi i m/12)\) en esa fila da la tabla real
anterior. Todos los valores son positivos. El menor de los dos valores
con \(\sqrt3\) es mayor que \(589609\), pues
\(658416>345420\sqrt3\), desigualdad que se comprueba elevando al
cuadrado sus dos miembros positivos. Los demás valores de la tabla son
mayores o iguales a \(589609\). El modo uniforme tiene valor
\((\sum K_i)^2=6263^2\); la desigualdad triangular de la transformada
de Fourier, usando \(K_i\geq0\), acota todos los demás por ese mismo
valor. El teorema espectral da \eqref{lib01:cotas-orbita}. Centrar
el registro elimina sólo el modo uniforme; los once restantes conservan
los valores no nulos de la tabla.
\end{proof}

El determinante, con el orden de columnas y la orientación de \(S\)
fijados aquí, es
\[
 \det O_K=5483119259214020183850397930008283625.
\]
La ausencia de un período propio del vector, por sí sola, no probaría la
invertibilidad: lo decisivo es que no se anula ninguno de sus modos de
Fourier. La condición numérica es aproximadamente \(8.15643\). Si se
normaliza \(K\) a norma uno, los extremos del Gram se dividen por
\(\|K\|^2=4251257\); la amplitud de sus componentes no debe
confundirse con una mejora intrínseca frente a toda referencia de igual
norma.

\begin{corollary}[Separación polar de forma y orientación]
\label{lib01:polar}
Sean
\[
 G_K=O_K^*O_K,\qquad P_K=G_K^{1/2},\qquad
 F_K=O_KG_K^{-1/2}.
\]
Entonces \(P_K\) es positivo definido, \(F_K\) es ortogonal y
\(O_K=F_KP_K\). Una realización isométrica \(J:V\to\mathcal Y\)
conserva \(G_K\) y \(P_K\), y transporta la parte isométrica a
\(JF_K\).
\end{corollary}
\begin{proof}
La positividad del Gram permite definir su raíz e inversa mediante el
teorema espectral. La identidad
\(F_K^*F_K=G_K^{-1/2}G_KG_K^{-1/2}=I\) prueba la afirmación.
Además, \((JO_K)^*(JO_K)=O_K^*O_K\).
\end{proof}

Esta polar se define a partir del Gram completo: no se obtiene sólo de la
lista de sus valores propios ni de la norma de \(K\). El símbolo
\(P_K\) de este corolario no es el proyector dimensional de rango diez
que se construirá más abajo. Esta separación es una consecuencia de la
referencia cíclica, no una identificación de todos sus lectores.

\subsection{Identificación de operadores y holonomía relativa}
\label{lib01:identificacion}

\begin{theorem}[Doce respuestas generales o una respuesta cíclica]
\label{lib01:operadores}
Para cualquier operador lineal \(H:V\to V\), sus doce respuestas
\(Y=(HK,HSK,\ldots,HS^{11}K)\) determinan
\[
 H=YO_K^{-1}.
\]
Si \(HS=SH\), una única respuesta vectorial completa \(y=HK\)
determina
\begin{equation}
 H=O_yO_K^{-1},\qquad
 H=\sum_{j=0}^{11}h_jS^j,\qquad h=O_K^{-1}y.
 \label{lib01:filtro}
\end{equation}
Con errores \(E\) en la matriz de respuestas o \(\delta y\) en el
vector, respectivamente,
\begin{align}
 \|\widehat H-H\|&\leq\frac{\|E\|}{\sqrt{589609}},
 \label{lib01:error-matriz}\\
 \|\widehat h-h\|&\leq\frac{\|\delta y\|}{\sqrt{589609}},\qquad
 \|\widehat H-H\|\leq\sqrt{\frac{12}{589609}}\,\|\delta y\|.
 \label{lib01:error-filtro}
\end{align}
\end{theorem}
\begin{proof}
La primera igualdad es \(Y=HO_K\). Si \(H\) conmuta con \(S\),
entonces \(HS^jK=S^jy\), de modo que \(Y=O_y\). La conmutación con
el ciclo determina todas las columnas de \(H\) desde una de ellas;
las doce potencias \(I,S,\ldots,S^{11}\) son independientes y generan
ese espacio de operadores. Por tanto \(y=O_Kh\), lo que da la
segunda fórmula. Las dos primeras cotas usan
\(\|O_K^{-1}\|=1/\sqrt{589609}\). La última usa
\(\|O_{\delta y}\|\leq\|O_{\delta y}\|_{\rm F}
=\sqrt{12}\|\delta y\|\).
\end{proof}

Una respuesta significa aquí un vector de doce componentes, no una
medición escalar. Como ejemplo, si \(H=2I+3S-S^5\), la respuesta
\(y=HK\) devuelve exactamente los coeficientes \(2,3,-1\) en sus
posiciones mediante \eqref{lib01:filtro}; el inversor necesita \(y,K,S\),
no los coeficientes del filtro. En cambio, el operador no nulo cuya primera
fila es \((543,-234,0,\ldots,0)\) y cuyas restantes filas son nulas
anula \(K\). Una sola respuesta no identifica un operador arbitrario.

La función de referencia admite una generalización exacta. En un ciclo de
\(n\) posiciones, \(H\mapsto Hg\) identifica todos los operadores
que conmutan con el ciclo si y sólo si
\(O_g=(g,Sg,\ldots,S^{n-1}g)\) es invertible. La suficiencia es la
prueba anterior. Si \(O_g\) es singular, un vector no nulo de su núcleo
da un polinomio de grado menor que \(n\) que anula \(g\); el operador
es no nulo porque las potencias del ciclo son independientes. El impulso
unitario también tiene órbita ortonormal y sirve de referencia. La
consecuencia específica para \(K\) es que el mismo objeto ya usado en
clausura e incidencia satisface esa condición, sin haber sido elegido de
nuevo para optimizar el ensayo. No se deduce una identificación de
dinámicas no lineales ni de estados ocultos arbitrarios; una aplicación a
una linealización necesita que ésta se haya definido y justificado.

\begin{corollary}[Recuperación de holonomía relativa desde el complemento]
\label{lib01:holonomia}
Sean \(H_B\) conocido e invertible y \(H_C\) otro transporte del
mismo tipo. El complemento
\[
 D_\gamma=\frac{\sqrt8}{9}(H_C-H_B)
\]
determina
\[
 H_B^{-1}H_C=I+\frac9{\sqrt8}H_B^{-1}D_\gamma.
\]
Si \(D_\gamma\) conmuta con \(S\), en particular si lo hacen ambos
transportes, basta \(y=D_\gamma K\) para recuperar
\begin{equation}
 H_B^{-1}H_C
 =I+\frac9{\sqrt8}H_B^{-1}O_yO_K^{-1}.
 \label{lib01:holonomia-una-respuesta}
\end{equation}
El error en esta reconstrucción está acotado por
\[
 \frac9{\sqrt8}\|H_B^{-1}\|
 \sqrt{\frac{12}{589609}}\,\|\delta y\|.
\]
\end{corollary}
\begin{proof}
Se despeja \(H_C\) en la definición del complemento y se aplica
\cref{lib01:operadores}. La cota resulta de la submultiplicatividad.
\end{proof}

Para \(H_B=S^3\) y \(H_C=S^4\), el dato racional
\(y/\sqrt8=(S^4-S^3)K/9\) permite reconstruir la holonomía relativa
\(S\) sin proporcionar \(H_C\) al inversor. Si \(H_B\) es
unitario, el factor de error es aproximadamente \(0.01435510\). Sin
conmutación se mantienen las doce respuestas generales, no una hipótesis
ficticia de simetría. La selección de un lazo físico y su representación
siguen perteneciendo a su dinámica generadora; este protocolo no las
sustituye.

\subsection{La memoria completa conserva la referencia y los observables}
\label{lib01:memoria-completa}

Sean \(C_n\) operadores unitarios de un espacio de Hilbert \(H\),
\[
 T_n=\frac{8I+C_n}{9},\qquad
 D_n=\frac{\sqrt8(C_n-I)}9,\qquad
 R_0=I,\quad R_{n+1}=T_nR_n.
\]
La escritura incluye secuencias no estacionarias. No exige conmutación
entre \(C_n\) de etapas distintas.

\begin{theorem}[Registro isométrico, terminal y transporte]
\label{lib01:isometria}
Para todo \(N\),
\[
 R_N^*R_N+\sum_{n=0}^{N-1}R_n^*D_n^*D_nR_n=I.
\]
Existe el límite fuerte positivo
\(A_\infty=\lim_N R_N^*R_N\), y
\begin{equation}
 \mathcal Vx=
 \bigl(A_\infty^{1/2}x,(D_nR_nx)_{n\geq0}\bigr)
 \quad\text{satisface}\quad\mathcal V^*\mathcal V=I.
 \label{lib01:V-infinita}
\end{equation}
Su imagen \(\mathcal M\) es cerrada y la inversa sobre ella es
\(\mathcal V^*\), de norma uno. En el caso finito se sustituye el
primer componente por \(R_Nx\) y se conservan los \(N\)
complementos. Para cualquier operador acotado \(A\),
\[
 A^{\mathcal V}=\mathcal VA\mathcal V^*\big|_{\mathcal M}
\]
preserva productos, adjuntos, norma y espectro en esa imagen, y
\[
 \langle\mathcal Vx,A^{\mathcal V}\mathcal Vy\rangle
 =\langle x,Ay\rangle.
\]
\end{theorem}
\begin{proof}
Expandir los dos términos de una etapa da
\(T_n^*T_n+D_n^*D_n=I\). Por tanto
\(R_n^*R_n-R_{n+1}^*R_{n+1}=R_n^*D_n^*D_nR_n\), y la suma
telescópica demuestra la identidad finita. La sucesión de operadores
positivos \(R_N^*R_N\) es decreciente y acotada inferiormente; su
límite fuerte existe. Tomando productos con \(x\), la suma de las
normas cuadradas de los complementos es
\(\|x\|^2-\langle x,A_\infty x\rangle\). Esto prueba que la
serie pertenece a la suma hilbertiana y que \(\mathcal V\) es una
isometría. Una isometría tiene imagen cerrada e inversa adjunta en ella.
Insertar \(\mathcal V^*\mathcal V=I\) en las composiciones demuestra
las propiedades del transporte. No se ha supuesto convergencia de
\(R_Nx\), ni que \(A_\infty\) sea un proyector.
\end{proof}

Al aplicar este resultado al portador del registro,
\begin{equation}
 (\mathcal VO_K)^*(\mathcal VO_K)=O_K^*O_K.
 \label{lib01:gram-memoria}
\end{equation}
Se conservan las doce direcciones y todas sus correlaciones, con las
mismas cotas de \cref{lib01:marco-ciclico}, a cualquier profundidad de
memoria completa. No se afirma que el continuo completo tenga dimensión
doce: se conserva esta referencia de doce dimensiones en su imagen.
Los operadores correctos allí son
\(S^{\mathcal V}=\mathcal VS\mathcal V^*\) y
\(H^{\mathcal V}=\mathcal VH\mathcal V^*\). No se presume que un
desplazamiento ambiente de las casillas de almacenamiento sea
\(S^{\mathcal V}\).

Si \(z=\mathcal Vy+\delta z\), la estimación
\[
 \widehat h=O_K^{-1}\mathcal V^*z
\]
mantiene la primera cota de \eqref{lib01:error-filtro} con
\(\|\delta z\|\) en lugar de \(\|\delta y\|\). La profundidad
no empeora esa constante porque se conserva la isometría y su terminal.
Una selección parcial de canales tiene otro Gram y requiere sus propias
cotas. Las potencias exteriores \(\bigwedge^p\mathcal V\) conservan
los determinantes de Gram; también se conservan trazas de rango finito y
medidas espectrales de los observables transportados. Las normas
cuadradas no se interpretan por ello automáticamente como probabilidades
de Born: tal lectura exige su realización posterior.

\subsection{Incidencia excepcional y dirección seleccionada}
\label{lib01:direccion}

Sea \(e=\mathbf1\), \(\Pi=I-ee^*/12=P_{11}\). En la carta
icosaédrica \(A_5/C_5\) ya construida, la representación ortogonal
\(\rho\) y el carácter tridimensional fijan
\begin{equation}
 P_3=\frac3{60}\sum_{g\in A_5}\chi_3(g^{-1})\rho(g),
 \quad P_3^*=P_3=P_3^2,\quad
 \operatorname{rank}P_3=3,\quad P_3e=0.
 \label{lib01:P3}
\end{equation}
Este proyector se recibe con su representación y sus marcas, no se elige
para adaptar una dirección al registro. Pongamos
\[
 k=\Pi K,\qquad u=P_3k,\qquad
 E_u=\frac{uu^*}{\|u\|^2},\qquad Q_K=\Pi-E_u=P_{10}(K).
\]
Las evaluaciones exactas son
\begin{align}
 \|k\|^2&=\frac{11789915}{12},&
 \|u\|^2&=\frac{6638585+2275584\sqrt5}{20},
 \label{lib01:normas-direccion}\\
 \eta_K:=\frac{\|u\|^2}{\|k\|^2}
 &=\frac{3(6638585+2275584\sqrt5)}{5\cdot11789915}
 \simeq0.5967954228259091.
 \label{lib01:eta}
\end{align}
En particular, \(u\ne0\). Como \(u^*k=\|u\|^2\),
\begin{equation}
 E_uk=u,\qquad Q_Kk=k-u,\qquad
 k=u+Q_Kk,\qquad \|k\|^2=\|u\|^2+\|Q_Kk\|^2.
 \label{lib01:particion-direccion}
\end{equation}
La reducción \(12\to11\to10\) elimina primero la dirección uniforme
y después una recta seleccionada. No elimina todo \(\operatorname{im}P_3\):
en \(\operatorname{im}Q_K\) permanecen las dos direcciones de ese
subespacio ortogonales a \(u\). Estas dimensiones son las de los
espacios algebraicos escritos; su identificación con dimensiones físicas
requiere sus mapas de realización.

\begin{proposition}[Transporte exacto de dirección y descriptor]
\label{lib01:transporte-direccion}
Para una isometría completa \(\mathcal V\), sea
\(\Pi^{\mathcal V}=\mathcal V\Pi\mathcal V^*\),
\(P_3^{\mathcal V}=\mathcal VP_3\mathcal V^*\) y
\(u^{\mathcal V}=\mathcal Vu\). En \(\mathcal M\),
\[
 Q_K^{\mathcal V}
 =\mathcal VQ_K\mathcal V^*
 =\Pi^{\mathcal V}
 -\frac{u^{\mathcal V}(u^{\mathcal V})^*}{\|u^{\mathcal V}\|^2},
 \quad \operatorname{rank}Q_K^{\mathcal V}=10,
 \quad \eta_K^{\mathcal V}=\eta_K.
\]
Además, \(\mathcal V^*u^{\mathcal V}=u\) y
\(\mathcal V^*Q_K^{\mathcal V}\mathcal V=Q_K\).
\end{proposition}
\begin{proof}
Se sustituyen las definiciones y \(\mathcal V^*\mathcal V=I\).
Las normas y el rango se conservan por la isometría. La identidad de la
imagen es \(\mathcal V\mathcal V^*\), no la identidad de un
almacenamiento ambiente mayor. Extender por cero fuera de la imagen no
añade direcciones físicas.
\end{proof}

La pantalla incidencial completa ofrece otra realización exacta. Si
\(\mathcal I\) es el operador de incidencia ya construido y
\(E_{\rm inc}=\mathcal I(e^\perp)\), entonces
\begin{equation}
 F:V\longrightarrow Y=\mathbb R\oplus E_{\rm inc},\quad
 Fx=\left(\mu(x),\frac{\mathcal I\Pi x}{6}\right),\quad
 \mu(x)=\frac{\mathbf1^*x}{12},\qquad
 \langle(\mu,y),(\nu,z)\rangle_Y=12\mu\nu+\langle y,z\rangle
 \label{lib01:incidencia-F}
\end{equation}
es una isometría sobreyectiva, pues
\(\Pi\mathcal I^*\mathcal I\Pi=36\Pi\). Para la incidencia
cruda de las 132 hexadas, la identidad completa es
\(\mathcal I^*\mathcal I=36I+30\mathbf1\mathbf1^*\); el
centrado elimina exactamente su componente uniforme. La imagen incidencial tiene
dimensión once; \(Y\) no es todo \(\mathbb R\oplus\mathbb R^{132}\).
Aplicar \(F\) término a término a la memoria conserva tanto el Gram
como los operadores anteriores y sus dominios. No se introduce por esta
composición una acción del Monstruo sobre \(K\).

\subsection{Estabilidad continua y necesidad de conservar correlaciones}
\label{lib01:ruido}

\begin{theorem}[Error de la dirección desde memoria completa]
\label{lib01:error-proyectores}
Sea \(z=\mathcal VK+\delta\), \(\|\delta\|\leq\varepsilon\),
y reconstruyamos
\(\widehat K=\mathcal V^*z\), \(\widehat k=\Pi\widehat K\),
\(\widehat u=P_3\widehat k\). Los tres errores son menores o iguales
a \(\varepsilon\). Si \(\varepsilon<\|u\|\),
\[
 \|Q_{\widehat K}-Q_K\|\leq\frac\varepsilon{\|u\|}.
\]
Si \(\varepsilon<\|k\|\),
\[
 |\eta_{\widehat K}-\eta_K|\leq\frac\varepsilon{\|k\|}.
\]
\end{theorem}
\begin{proof}
El adjunto de la isometría y los proyectores tienen norma uno. El ruido
ortogonal a \(\mathcal M\) desaparece al aplicar \(\mathcal V^*\).
Para dos rectas no nulas, la norma de la diferencia de sus proyectores es
el seno del ángulo entre ellas. La distancia de \(u\) a la recta de
\(\widehat u\) no supera \(\|u-\widehat u\|\), por lo que ese
seno es menor o igual que \(\varepsilon/\|u\|\). La hipótesis
impide \(\widehat u=0\).

Para la segunda afirmación se normalizan \(k\) y \(\widehat k\).
La diferencia de sus proyectores de rango uno tiene, en su plano, valores
propios \(\pm\sin\theta\). Al emparejarla con \(0\leq P_3\leq I\)
el valor absoluto no supera \(\sin\theta\): en una base de vectores
propios se obtiene \(\sin\theta(a-b)\), con \(a,b\in[0,1]\).
El mismo argumento de distancia da
\(\sin\theta\leq\varepsilon/\|k\|\). Si las rectas coinciden,
la diferencia de los cocientes es cero.
\end{proof}

Aquí \(\|u\|\simeq765.733\), por lo que el factor para el proyector
es aproximadamente \(0.00130594\). En el lector de sólo dos memorias,
la inversa tiene norma \(9/4\) y el factor correspondiente es
\(9/(4\|u\|)\simeq0.00293836\). Son datos observados distintos;
la diferencia no es una mejora gratuita del mismo canal. La corrección
entera de \cref{lib02:radios-memoria} permitirá recuperar primero el
registro sin error dentro de su radio, en lugar de propagar un error
continuo hasta el proyector.

\begin{proposition}[Los bloques cruzados forman parte del observable]
\label{lib01:correlaciones}
Escribamos \(\mathcal Vx=(W_ax)_a\), incluyendo el terminal. Entonces
el bloque \((a,b)\) de \(A^{\mathcal V}\) es
\[
 (A^{\mathcal V})_{ab}=W_aAW_b^*.
\]
Su forma cuadrática se calcula como límite de truncamientos finitos
\[
 \langle y,A^{\mathcal V}y\rangle
 =\lim_F\sum_{a,b\in F}\langle y_a,W_aAW_b^*y_b\rangle.
\]
La conservación general no se obtiene descartando todos los términos
\(a\ne b\).
\end{proposition}
\begin{proof}
La fórmula de bloques resulta de multiplicar \(\mathcal VA\mathcal V^*\).
Los proyectores de truncamiento convergen fuertemente a la identidad y
\(A^{\mathcal V}\) es acotado, lo que prueba el límite sin afirmar
convergencia absoluta de una serie doble arbitraria. Para un
contraejemplo, tome \(H=\mathbb R\), \(C=-I\),
\(T=7/9\), \(D=-2\sqrt8/9\), \(\mathcal Vx=(Tx,Dx)\) y
\(A=I\). La parte diagonal, evaluada en \(\mathcal Vx\), vale
\[
 (T^4+D^4)x^2=\frac{3425}{6561}x^2.
\]
Los términos cruzados aportan \(3136x^2/6561\), y la suma es
\(x^2\), como exige la isometría.
\end{proof}

\subsection{Corrección discreta antes de la selección incidencial}
\label{lib01:correccion-incidencia}

Los lectores canónicos de dos memorias son
\[
 T_j=\frac{8I+S^j}{9},\qquad
 D_j=\frac{\sqrt8(S^j-I)}9,\qquad
 Mx=(D_3x,D_4T_3x).
\]
La segunda memoria actúa después de la primera compresión; no se reinyecta
en este protocolo el primer complemento. El teorema general de
\cref{lib02:radios-memoria} demostrará, sin utilizar el registro como
objetivo, los radios estrictos
\[
 r_0=\frac{2\sqrt{146}}{81},\qquad
 r_q=\frac{4\sqrt{73}}{81}.
\]
El primer radio recupera \(\Pi K\) a partir de \(MK\), módulo la
traslación entera uniforme. El segundo recupera \(K\) íntegro cuando
se conserva además \(q=\mathbf1^*K\). Si se observa una carga real
con error menor que \(1/2\), su redondeo recupera antes el canal entero
de carga. Ninguno de estos procedimientos necesita el terminal
\(T_4T_3K\).

La selección de bloque alto en la carta marcada es
\begin{equation}
 B_K=\{i:K_i\geq729\}=\{4,6,9,10\}.
 \label{lib01:umbral}
\end{equation}
La cuarta componente está exactamente en el umbral. Esta circunstancia
separa la estabilidad continua de la discreta.

\begin{proposition}[Margen continuo nulo y restitución entera]
\label{lib01:umbral-estable}
El selector directo de \eqref{lib01:umbral}, aplicado a la reconstrucción
real de las memorias, no tiene radio euclídeo positivo de estabilidad en
\(MK\), incluso con carga exacta. En cambio, la decodificación entera
previa con error menor que \(r_q\) restituye exactamente \(B_K\)
sin modificar el umbral.
\end{proposition}
\begin{proof}
Sea \(L\) la inversa de mínimos cuadrados, que satisface
\(LM=\Pi\). Para \(t>0\), ponga
\[
 e_t=-tM\Pi e_4,
 \qquad \widetilde K_t=\frac q{12}\mathbf1+L(MK+e_t)
 =K-t\Pi e_4.
\]
Se tiene \(\|e_t\|\to0\), pero la cuarta componente es
\(729-11t/12<729\). Para \(t\) pequeño sólo cambia esa pertenencia.
La segunda afirmación sigue de recuperar primero \(K\) exactamente
mediante \cref{lib02:radios-memoria} y aplicar después el mismo selector.
\end{proof}

Si los restantes datos marcados \(P_\pi,H_e,H_\varphi,H_{\rm APP}\)
se conservan exactos o se transportan por su ley, la unicidad de la
construcción excepcional devuelve la misma bandera
\(B_K\subset H^-_\alpha\) y el mismo origen marcado. Corregir el
registro no equivale a corregir errores independientes de esas marcas.
El testigo
\[
 K'=K-e_4+e_6,
 \qquad Q(K')=6263,\qquad B_{K'}=\{6,9,10\}
\]
pertenece también a \(\{0,\ldots,999\}^{12}\) y cumple
\(\|MK'-MK\|^2=4672/6561\). El punto medio de ambas memorias está
a distancia \(r_q\) de cada una. Prueba que el radio estricto es
óptimo uniformemente en el dominio algebraico de carga fija, no que
\(K'\) tenga una historia terminal HMT construida.

Sin carga, \(M(K-\mathbf1)=MK\) y las selecciones altas son distintas.
Las dos memorias pueden recuperar exactamente \(Q_K\) y \(\eta_K\),
que ignoran la media, sin recuperar ese selector absoluto. Para el
proyector sigue siendo necesaria la condición \(P_3\Pi K\ne0\);
la corrección entera no suprime esa condición de definición. Un ensayo
finito concreto usa \(z=MK/\sqrt8\) y añade \(1/10\) a su primera
componente. El error en la norma de \(M\) es \(\sqrt8/10<r_q\).
El algoritmo de suelo y techo de \cref{lib02:decodificador} deja 924
candidatos de carga 6263 y recupera \(K\) y \(B_K\) exactamente.

Aplicando la isometría \eqref{lib01:incidencia-F} se obtiene
\(M_{\rm inc}=(F\oplus F)M\). Todas las distancias entre palabras
y las normas de ruido se conservan, de modo que los mismos radios y
decodificadores valen en \(Y\oplus Y\). Ésta es la conexión
cuantitativa entre memoria e incidencia; no depende de atribuir al
desplazamiento \(S\) una simetría excepcional adicional.

\subsection{Covarianza, refinamiento y carga variacional}
\label{lib01:covarianza}

Si \(G\) es ortogonal y se transportan conjuntamente
\(K'=GK\), \(\Pi'=G\Pi G^*\) y \(P_3'=GP_3G^*\), entonces
\[
 u'=Gu,\qquad Q_{K'}'=GQ_KG^*,\qquad \eta_{K'}'=\eta_K.
\]
En una carta fija, \(K\mapsto aK+b\mathbf1\), \(a\ne0\), deja
invariantes \(Q_K\) y \(\eta_K\); no deja invariante todo el objeto.
Por ejemplo, \(\kappa(K+\mathbf1)-\kappa(K)=1/999\), y la lectura
\(\alpha_A\) cambiaría en \(-1/999\) si se mantuviesen fijas las
regiones. Es un contraejemplo algebraico a identificar el descriptor con
todo el registro, no una afirmación de admisibilidad de esa actualización
como historia HMT.

Tampoco se puede mover sólo el registro y mantener arbitrariamente el
proyector. La fórmula finita \eqref{lib01:P3} da exactamente
\[
 \|P_3S^3K\|^2=\frac{369035}{4}-\frac{281861\sqrt5}{10},
 \qquad
 \|P_3S^4K\|^2=\frac{1327717}{4}+\frac{694618\sqrt5}{5}.
\]
Son distintas de \eqref{lib01:normas-direccion}; en particular,
\([P_3,S^3]\) y \([P_3,S^4]\) no se anulan. Al transportar
\(P_3\) por la misma conjugación sí se conserva \(\eta_K\).

La decodificación también es covariante. Para una permutación \(R\),
\[
 S'=RSR^{-1},\qquad M'=(R\oplus R)MR^{-1},\qquad
 M'R=(R\oplus R)M.
\]
La red, la carga y las distancias se conservan. En el régimen de unicidad,
la salida del decodificador se transforma por \(R\); fuera de él se
transporta el conjunto de minimizadores, sin elegir un empate de forma
arbitraria. Un marco ortogonal general exige transportar también
\(\mathbb Z^{12}\) a \(R\mathbb Z^{12}\) y la dirección uniforme
a \(R\mathbf1\). Redondear en la red antigua después de cambiar la
carta no representa el mismo problema.

\begin{proposition}[Naturalidad bajo refinamiento con marcos locales]
\label{lib01:refinamiento}
Sea \(z\) un cilindro con hijos \(z'\), pesos positivos compatibles
\(\sum_{\tau z'=z}\mu(z')=\mu(z)\), y marcos ortogonales \(G_z\).
El refinamiento
\[
 \mathcal R_n(e_z\otimes x)=
 \sum_{\tau z'=z}\sqrt{\frac{\mu(z')}{\mu(z)}}
 e_{z'}\otimes G_{z'}G_z^*x
\]
es isométrico. Para el campo de proyectores
\(Q_z=G_zQ_KG_z^*\) y su operador por bloques \(Q_n\),
\[
 Q_{n+1}\mathcal R_n=\mathcal R_nQ_n,
 \qquad \mathcal R_n^*Q_{n+1}\mathcal R_n=Q_n.
\]
Las identidades se componen entre niveles y pasan al límite inductivo.
\end{proposition}
\begin{proof}
Los hijos tienen soportes ortogonales y los pesos suman uno; los cambios
de marco son ortogonales. En cada hijo,
\(Q_{z'}G_{z'}G_z^*=G_{z'}G_z^*Q_z\), que demuestra el primer
cuadrado. Multiplicar por el adjunto del refinamiento da el segundo.
La compatibilidad por composición define el operador en el límite de
las inclusiones isométricas. El mismo argumento vale para \(E_u\).
\end{proof}

El enunciado conserva la publicación de un prefijo bajo refinamientos
compatibles. Nuevos eventos pueden producir otro registro; no se deduce
una constancia temporal de \(K\). Los rangos de campos sobre fibras
tensoriales son los de esas fibras, no un recuento de dimensiones físicas.

\begin{proposition}[Una carga de transporte con acción explícita]
\label{lib01:carga-noether}
Sean \(U_n:E_n\to E_{n+1}\) transportes invertibles. La acción
auxiliar discreta con variables cotangentes es
\[
 \mathcal A(q,p)=\sum_n p_{n+1}^*(q_{n+1}-U_nq_n).
\]
Las variaciones de soporte finito dan
\(q_{n+1}=U_nq_n\) y \(p_n=U_n^*p_{n+1}\). Si
\(G_0=I\), \(G_{n+1}=U_nG_n\), \(A_0=Q_K\) y
\(A_n=G_nA_0G_n^{-1}\), la cantidad
\[
 J_n=p_n^*A_nq_n
\]
es independiente de \(n\). La acción es invariante bajo el grupo
continuo
\(q_n\mapsto\exp(tA_n)q_n\),
\(p_n\mapsto\exp(-tA_n^*)p_n\).
\end{proposition}
\begin{proof}
Las ecuaciones indicadas resultan al variar \(p_{n+1}\) y \(q_n\).
La definición de \(A_n\) implica
\(A_{n+1}U_n=U_nA_n\). Por tanto
\[
 p_{n+1}^*A_{n+1}q_{n+1}
 =p_{n+1}^*U_nA_nq_n=p_n^*A_nq_n.
\]
El mismo entrelazamiento vale para las exponenciales y hace que cada
término de la acción sea invariante. No se requiere que \(A_n\)
sea antisimétrico: la transformación del covector es la contragrediente.
\end{proof}

Para transportes ortogonales puede elegirse \(p_n=q_n\); entonces
\(J_n=\|A_nq_n\|^2\), y el complemento proporciona la carga
restante. En una realización simpléctica sin forma positiva, se conserva
el emparejamiento escrito y no se lo llama energía positiva. Si se
archivan las fibras mediante isometrías \(\mathcal V_n\), las
definiciones
\[
 \widehat U_n=\mathcal V_{n+1}U_n\mathcal V_n^*,\qquad
 \widehat A_n=\mathcal V_nA_n\mathcal V_n^*,\qquad
 \widehat q_n=\mathcal V_nq_n,
\]
con \(\widehat p_n(y)=p_n(\mathcal V_n^*y)\), conservan el
entrelazamiento y la misma carga en las imágenes. Esta composición
proporciona una acción y una simetría continuas concretas para una ley de
conservación; no atribuye corrientes de Noether automáticas a un grupo
finito, ni identifica la acción auxiliar con toda acción física.

\subsection{Codificación, identificación operatoria y estabilidad de la incidencia}
\label{lib01:alcance}

La aportación reunida es una cadena operativa: el registro generado tiene
una codificación reversible; su órbita es una referencia completa;
terminal y complementos conservan esa referencia y los observables
transportados; la separación entera de dos memorias permite restituir
incidencia y dirección aun cuando el selector real tiene margen cero.
Las pruebas de linealidad, de proyección y de redes se utilizan como
lenguaje de realización posterior, sin cambiar la producción del objeto.

Cada hipótesis tiene un testigo que impide ampliarla inadvertidamente.
Una respuesta aislada no identifica un operador que no conmuta con el
ciclo; sólo la memoria completa conserva automáticamente el Gram original;
los términos cruzados no pueden omitirse; la selección absoluta exige
carga; el proyector requiere dirección no nula; las cotas de ruido
dependen de la normalización; una transformación de carta exige transportar
forma, red y lectores. Una perturbación de lectura no se convierte por
decreto en una historia admisible. Estas delimitaciones preservan las
funciones de \(K\), sin reducirlo a doce números desconectados ni
promover una propiedad canónica a invariancia universal.
